\section{Practical Computation and Approximation of \texorpdfstring{$\Sigma_2$}{Sigma-2} Objects}
\label{sec:sigma2-practical}

A natural objection is that $\Sigma_2$ objects appear computationally intractable: specifying
$\mathcal{F}(r)=\Pi^{-1}(r)$ and the relations $\{\mathcal{A}_I(r)\}$ exactly may be as hard as
working with $\mathcal{H}$ itself. The intent of $\Sigma_2$ is not to demand full reconstruction of
$\mathcal{H}$, but to identify the \emph{structurally necessary} information for counterfactual
prediction and to provide a target for approximation schemes.

\subsection{Finite parameterizations}

In many physically relevant cases, histories within $\mathcal{F}(r)$ differ only through a
finite set of parameters (e.g., a low-rank family of initial correlations, a small number of
bath modes, or a finite memory length). In such settings, one may represent $\mathcal{F}(r)$
implicitly as a parameterized family
\begin{equation}
\mathcal{F}(r) \approx \{ h(\theta) : \theta \in \Theta_r \},
\end{equation}
and approximate admissibility relations by pushforward of $\sim_I$ through this parameterization.

\subsection{Sample-based approximations}

When explicit parameterization is unavailable, $\Sigma_2$ admits sample-based approximations. One may
maintain a finite ensemble $\{h_j\}_{j=1}^N \subset \mathcal{F}(r)$ and approximate continuation
sets by evaluating which reduced outcomes are reachable under interventions for sampled elements.
This yields empirical approximations of the form
\begin{equation}
\widehat{\mathrm{Cont}}(r;\mathcal{I})
=
\{(I,\Pi(h_j')) : h_j \sim_I h_j' \text{ for some } j \}.
\end{equation}
Such approximations are directly relevant to robust control and verification problems, where the
objective is often to guarantee behavior over a set of possible hidden correlations rather than
to infer a unique hidden state.

\subsection{Coarse-graining by continuation equivalence}

Exact fibers are often too large, but $\Sigma_2$ suggests a principled coarse-graining: partition
$\mathcal{F}(r)$ into equivalence classes whose elements share identical continuation structure
under a restricted intervention class. Define an equivalence relation on $\mathcal{F}(r)$ by
\begin{equation}
h_1 \equiv_{\mathcal{I}} h_2
\iff
\forall I\in\mathcal{I}:\;
\{r':(h_1,r')\in\mathcal{A}_I(r)\}
=
\{r':(h_2,r')\in\mathcal{A}_I(r)\}.
\end{equation}
A $\Sigma_2$ description may then be represented by the quotient space $\mathcal{F}(r)/{\equiv_{\mathcal{I}}}$
together with the induced admissibility structure. This yields the \emph{coarsest representation
that preserves counterfactual prediction for the intervention class of interest}.

\subsection{Tractability and Structural Compression}
\label{sec:tractability-compression}

It is important to distinguish a $\Sigma_2$ description from a full
system--environment simulation. Although the fiber
$\mathcal{F}(r)=\Pi^{-1}(r)$ is defined as a set of histories, $\Sigma_2$
does not require enumerating or evolving those histories individually.
Rather, it tracks equivalence classes of histories that share identical
continuation structure with respect to the intervention class
$\mathcal{I}$.

In many physical regimes, the admissibility relations
$\{\mathcal{A}_I(r)\}_{I\in\mathcal{I}}$ collapse large families of
microscopically distinct histories into a low-dimensional
parameterization. In such cases, $\Sigma_2$ is strictly coarser than the
full history space $\mathcal{H}$ while remaining sufficient for
counterfactual prediction.

Examples of compressible fibers include:
\begin{itemize}
  \item finite-dimensional system--environment correlations, where the
  fiber is parameterized by a small set of correlation coefficients;
  \item bounded-memory reservoirs, where history dependence truncates
  after a finite time $\tau$;
  \item coarse-grained control-relevant observables, for which
  microscopically distinct states are indistinguishable to the available
  control bandwidth;
  \item experimentally indistinguishable preparation classes.
\end{itemize}

\subsection{The role of \texorpdfstring{$\Sigma_2$}{Sigma-2} in practice}

Even when full $\Sigma_2$ specification is impractical, the framework provides concrete benefits:
\begin{enumerate}
\item It supplies a diagnostic for when $\Sigma_1$ modeling is structurally incapable of supporting the
      counterfactual questions one intends to ask.
\item It defines the correct target object for approximation: not a single state or tensor, but a
      continuation-structured family indexed by hidden correlation classes.
\item It clarifies what must be preserved by reduced models and coarse-graining procedures to remain
      operationally faithful with respect to a chosen intervention class.
\end{enumerate}